% GENERATED FILE. Edit canonical lessons, then run npm run generate:books.
% canonical-source-hash: fnv1a64:add539dd90e7f724
% canonical-lessons: KA-C338-vadanti, KA-C338-pasarisu, KA-C338-uhapoha, KA-C338-satya, KA-C338-alisu-listen

\chapter{Rumours in the Town}
\label{ch:ka-rumours-in-the-town}
\hlchaptermodality{338}

\begin{chapteropening}
\textbf{By the end of this chapter:} \emph{I can say a rumour, to spread far, speculation, the truth and to listen closely.}

Complete the last lesson of chapter 338: I can say a rumour, to spread far, speculation, the truth and to listen closely.
\end{chapteropening}

\section[\texorpdfstring{\kn{ವದಂತಿ} (vadanti)}{ವದಂತಿ (vadanti)}]{\kn{ವದಂತಿ} (vadanti) — a rumour}

\label{lesson:KA-C338-vadanti}



\begin{quote}
Before the new one: say the Kannada for a tribute, then the Kannada for to comfort.
\end{quote}

\subsection*{You'll want to know: \kn{ವದಂತಿ}}
\textbf{\kn{ವದಂತಿ}} — \emph{vadanti} — \textquotedblleft{}a rumour\textquotedblright{}.

\begin{quote}
\kn{ಶಾಲೆ} \kn{ಮುಚ್ಚುತ್ತಾರೆ} \kn{ಎಂಬ} \kn{ವದಂತಿ} \kn{ಊರಿನಲ್ಲಿ} \kn{ಇತ್ತು}.
\end{quote}

\begin{quote}
\emph{śāle muccuttāre emba vadanti ūrinalli ittu.}
\end{quote}

\begin{quote}
There was a rumour in town that the school would close.
\end{quote}

Plural: \kn{ವದಂತಿಗಳು}, \emph{vadantigaḷu}.

A rumour \textquotedblleft{}spreads\textquotedblright{}: \kn{ವದಂತಿ} \kn{ಹಬ್ಬಿತು}. \kn{ಎಂಬ}, \textquotedblleft{}called, that says\textquotedblright{}, links it to its content.

\begin{cousinweb}
\textbf{\kn{ವದಂತಿ}} (\emph{vadanti}) is borrowed from Sanskrit \emph{vadanti}, a noun that Monier-Williams gives beside \emph{kiṃvadantī}, \textquotedblleft{}the common saying, a rumour\textquotedblright{}, literally \textquotedblleft{}what do they say\textquotedblright{}, from \emph{vad}, \textquotedblleft{}to speak\textquotedblright{}.
\end{cousinweb}

\begin{grammarlens}[title={What you've built}]
One more word for telling what happened.
\end{grammarlens}

\subsection*{Your turn}
\begin{itemize}
\raggedright
  \item \emph{Say it:} \emph{vadanti}
  \item \emph{Say it:} \emph{vadanti}, once more
  \item \emph{Say it:} \emph{śāle muccuttāre emba vadanti ūrinalli ittu.} — There was a rumour in town that the school would close.
  \item \emph{Recall:} say the Kannada for a tribute, then the Kannada for to comfort, then say \emph{vadanti} again
\end{itemize}

\begin{tcolorbox}[breakable,colback=teal!4,colframe=teal!35!black,title={Before you move on}]
What does \kn{ವದಂತಿ} mean? (\textquotedblleft{}A rumour\textquotedblright{}.) Say it once more.
\end{tcolorbox}
\section[\texorpdfstring{\kn{ಪಸರಿಸು} (pasarisu)}{ಪಸರಿಸು (pasarisu)}]{\kn{ಪಸರಿಸು} (pasarisu) — to spread far}

\label{lesson:KA-C338-pasarisu}



\begin{quote}
Before the new one: say the Kannada for to comfort, then the Kannada for a rumour.
\end{quote}

\subsection*{You'll want to know: \kn{ಪಸರಿಸು}}
\textbf{\kn{ಪಸರಿಸು}} — \emph{pasarisu} — \textquotedblleft{}to spread far\textquotedblright{}.

\begin{quote}
\kn{ಆ} \kn{ಸುದ್ದಿ} \kn{ನಗರದ} \kn{ತುಂಬಾ} \kn{ಪಸರಿಸಿತು}.
\end{quote}

\begin{quote}
\emph{ā suddi nagarada tumbā pasarisitu.}
\end{quote}

\begin{quote}
That news spread all over the city.
\end{quote}

Past (it …): \kn{ಪಸರಿಸಿತು}, \emph{pasarisitu}.

News, fame and fragrance spread far: \kn{ಸುಗಂಧ} \kn{ಪಸರಿಸಿತು}, a fragrance spread.

\begin{cousinweb}
\textbf{\kn{ಪಸರಿಸು}} (\emph{pasarisu}) is built on Sanskrit \emph{prasara}, \textquotedblleft{}advance, extension\textquotedblright{}, in its Prakrit shape \emph{pasara} (CDIAL 8824), with the Kannada verb ending -\emph{isu}.
\end{cousinweb}

\begin{grammarlens}[title={What you've built}]
One more word for telling what happened.
\end{grammarlens}

\subsection*{Your turn}
\begin{itemize}
\raggedright
  \item \emph{Say it:} \emph{pasarisu}
  \item \emph{Say it:} \emph{pasarisu}, once more
  \item \emph{Say it:} \emph{ā suddi nagarada tumbā pasarisitu.} — That news spread all over the city.
  \item \emph{Recall:} say the Kannada for to comfort, then the Kannada for a rumour, then say \emph{pasarisu} again
\end{itemize}

\begin{tcolorbox}[breakable,colback=teal!4,colframe=teal!35!black,title={Before you move on}]
What does \kn{ಪಸರಿಸು} mean? (\textquotedblleft{}To spread far\textquotedblright{}.) Say it once more.
\end{tcolorbox}
\section[\texorpdfstring{\kn{ಊಹಾಪೋಹ} (ūhāpōha)}{ಊಹಾಪೋಹ (ūhāpōha)}]{\kn{ಊಹಾಪೋಹ} (ūhāpōha) — speculation}

\label{lesson:KA-C338-uhapoha}



\begin{quote}
Before the new one: say the Kannada for a rumour, then the Kannada for to spread far.
\end{quote}

\subsection*{You'll want to know: \kn{ಊಹಾಪೋಹ}}
\textbf{\kn{ಊಹಾಪೋಹ}} — \emph{ūhāpōha} — \textquotedblleft{}speculation\textquotedblright{}.

\begin{quote}
\kn{ಆ} \kn{ಘಟನೆಯ} \kn{ಬಗ್ಗೆ} \kn{ತುಂಬಾ} \kn{ಊಹಾಪೋಹ} \kn{ಇತ್ತು}.
\end{quote}

\begin{quote}
\emph{ā ghaṭaneya bagge tumbā ūhāpōha ittu.}
\end{quote}

\begin{quote}
There was a lot of speculation about that incident.
\end{quote}

\kn{ಊಹೆ}, a guess, and \kn{ಅಪೋಹ}, a denial: guesses for and against.

\begin{cousinweb}
\textbf{\kn{ಊಹಾಪೋಹ}} (\emph{ūhāpōha}) is borrowed from Sanskrit, a compound of \emph{ūha}, \textquotedblleft{}consideration\textquotedblright{} (CDIAL 2443), and \emph{apoha}, \textquotedblleft{}reasoning, denying\textquotedblright{} (Monier-Williams).
\end{cousinweb}

\begin{grammarlens}[title={What you've built}]
One more word for telling what happened.
\end{grammarlens}

\subsection*{Your turn}
\begin{itemize}
\raggedright
  \item \emph{Say it:} \emph{ūhāpōha}
  \item \emph{Say it:} \emph{ūhāpōha}, once more
  \item \emph{Say it:} \emph{ā ghaṭaneya bagge tumbā ūhāpōha ittu.} — There was a lot of speculation about that incident.
  \item \emph{Recall:} say the Kannada for a rumour, then the Kannada for to spread far, then say \emph{ūhāpōha} again
\end{itemize}

\begin{tcolorbox}[breakable,colback=teal!4,colframe=teal!35!black,title={Before you move on}]
What does \kn{ಊಹಾಪೋಹ} mean? (\textquotedblleft{}Speculation\textquotedblright{}.) Say it once more.
\end{tcolorbox}
\section[\texorpdfstring{\kn{ಸತ್ಯ} (satya)}{ಸತ್ಯ (satya)}]{\kn{ಸತ್ಯ} (satya) — the truth; true}

\label{lesson:KA-C338-satya}



\begin{quote}
Before the new one: say the Kannada for to spread far, then the Kannada for speculation.
\end{quote}

\subsection*{You'll want to know: \kn{ಸತ್ಯ}}
\textbf{\kn{ಸತ್ಯ}} — \emph{satya} — \textquotedblleft{}the truth; true\textquotedblright{}.

\begin{quote}
\kn{ಕೊನೆಗೆ} \kn{ಸತ್ಯ} \kn{ಎಲ್ಲರಿಗೂ} \kn{ಗೊತ್ತಾಯಿತು}.
\end{quote}

\begin{quote}
\emph{konege satya ellarigū gottāyitu.}
\end{quote}

\begin{quote}
In the end everyone came to know the truth.
\end{quote}

\kn{ನಿಜ} is the everyday word for \textquotedblleft{}true\textquotedblright{}; \kn{ಸತ್ಯ} is weightier: \kn{ಸತ್ಯ} \kn{ಹೇಳು}, to tell the truth.

\begin{cousinweb}
\textbf{\kn{ಸತ್ಯ}} (\emph{satya}) is borrowed from Sanskrit \emph{satya}, \textquotedblleft{}true\textquotedblright{}, and as a noun \textquotedblleft{}truth\textquotedblright{} (CDIAL 13112).
\end{cousinweb}

\begin{grammarlens}[title={What you've built}]
One more word for telling what happened.
\end{grammarlens}

\subsection*{Your turn}
\begin{itemize}
\raggedright
  \item \emph{Say it:} \emph{satya}
  \item \emph{Say it:} \emph{satya}, once more
  \item \emph{Say it:} \emph{konege satya ellarigū gottāyitu.} — In the end everyone came to know the truth.
  \item \emph{Recall:} say the Kannada for to spread far, then the Kannada for speculation, then say \emph{satya} again
\end{itemize}

\begin{tcolorbox}[breakable,colback=teal!4,colframe=teal!35!black,title={Before you move on}]
What does \kn{ಸತ್ಯ} mean? (\textquotedblleft{}The truth; true\textquotedblright{}.) Say it once more.
\end{tcolorbox}
\section[\texorpdfstring{\kn{ಆಲಿಸು} (ālisu)}{ಆಲಿಸು (ālisu)}]{\kn{ಆಲಿಸು} (ālisu) — to listen closely}

\label{lesson:KA-C338-alisu-listen}



\begin{quote}
Before the new one: say the Kannada for speculation, then the Kannada for the truth.
\end{quote}

\subsection*{You'll want to know: \kn{ಆಲಿಸು}}
\textbf{\kn{ಆಲಿಸು}} — \emph{ālisu} — \textquotedblleft{}to listen closely\textquotedblright{}.

\begin{quote}
\kn{ಎಲ್ಲರೂ} \kn{ಅಜ್ಜನ} \kn{ಕಥೆಯನ್ನು} \kn{ಗಮನವಿಟ್ಟು} \kn{ಆಲಿಸಿದರು}.
\end{quote}

\begin{quote}
\emph{ellarū ajjana katheyannu gamanaviṭṭu ālisidaru.}
\end{quote}

\begin{quote}
Everyone listened closely to Grandfather's story.
\end{quote}

Past (I …): \kn{ಆಲಿಸಿದೆ}, \emph{āliside}.

A more careful listening than \kn{ಕೇಳು}; what one listens to takes -\kn{ಅನ್ನು}.

\begin{cousinweb}
\textbf{\kn{ಆಲಿಸು}} (\emph{ālisu}) is native Dravidian: DEDR 383 lists Kannada \emph{ālisu}, \emph{ālaysu}, \textquotedblleft{}to listen, to attend to\textquotedblright{}.
\end{cousinweb}

\begin{grammarlens}[title={What you've built}]
One more word for telling what happened.
\end{grammarlens}

\subsection*{Your turn}
\begin{itemize}
\raggedright
  \item \emph{Say it:} \emph{ālisu}
  \item \emph{Say it:} \emph{ālisu}, once more
  \item \emph{Say it:} \emph{ellarū ajjana katheyannu gamanaviṭṭu ālisidaru.} — Everyone listened closely to Grandfather's story.
  \item \emph{Recall:} say the Kannada for speculation, then the Kannada for the truth, then say \emph{ālisu} again
\end{itemize}

\begin{tcolorbox}[breakable,colback=teal!4,colframe=teal!35!black,title={Before you move on}]
What does \kn{ಆಲಿಸು} mean? (\textquotedblleft{}To listen closely\textquotedblright{}.) Say it once more.
\end{tcolorbox}
